% GENERATED FILE. Edit canonical lessons, then run npm run generate:books.
% canonical-source-hash: fnv1a64:b26de80e2a0c05aa
% canonical-lessons: GU-C486-arogya, GU-C486-svasth, GU-C486-sphurti, GU-C486-khorak, GU-C486-sukavu, GU-R486-review-479-486

\chapter{Staying Well}
\label{ch:gu-staying-well}
\hlchaptermodality{486}

\begin{chapteropening}
\textbf{By the end of this chapter:} \emph{I can say what keeps me healthy.}

Complete the last lesson of chapter 486: I can say what keeps me healthy.
\end{chapteropening}

\section[\texorpdfstring{\gu{આરોગ્ય} (ārogya)}{આરોગ્ય (ārogya)}]{\gu{આરોગ્ય} (ārogya) — good health}

\label{lesson:GU-C486-arogya}



\begin{quote}
Before the new one: say the Gujarati for a dose, then the Gujarati for a diet restriction.
\end{quote}

\subsection*{You'll want to know: \gu{આરોગ્ય}}
\textbf{\gu{આરોગ્ય}} — \emph{ārogya} — \textquotedblleft{}good health\textquotedblright{}.

\begin{quote}
\gu{રોજ} \gu{ચાલવું} \gu{આરોગ્ય} \gu{માટે} \gu{સારું} \gu{છે}.
\end{quote}

\begin{quote}
\emph{roj cālvũ ārogya māṭe sārũ che.}
\end{quote}

\begin{quote}
Walking every day is good for your health.
\end{quote}

A neuter noun.

\begin{cousinweb}
\textbf{\gu{આરોગ્ય}} (\emph{ārogya}) is a tatsama, Sanskrit \emph{ārogya}, \textquotedblleft{}freedom from disease, health\textquotedblright{}, as Monier-Williams records it; CDIAL lists an older spoken form, \emph{ārog} (CDIAL 1330).
\end{cousinweb}

\begin{grammarlens}[title={What you've built}]
One more word for describing what you lived and felt.
\end{grammarlens}

\subsection*{Your turn}
\begin{itemize}
\raggedright
  \item \emph{Say it:} \emph{ārogya}
  \item \emph{Say it:} \emph{ārogya}, once more
  \item \emph{Say it:} \emph{roj cālvũ ārogya māṭe sārũ che.} — Walking every day is good for your health.
  \item \emph{Recall:} say the Gujarati for a dose, then the Gujarati for a diet restriction, then say \emph{ārogya} again
\end{itemize}

\begin{tcolorbox}[breakable,colback=teal!4,colframe=teal!35!black,title={Before you move on}]
What does \gu{આરોગ્ય} mean? (\textquotedblleft{}Good health\textquotedblright{}.) Say it once more.
\end{tcolorbox}
\section[\texorpdfstring{\gu{સ્વસ્થ} (svasth)}{સ્વસ્થ (svasth)}]{\gu{સ્વસ્થ} (svasth) — in good health; well}

\label{lesson:GU-C486-svasth}



\begin{quote}
Before the new one: say the Gujarati for a diet restriction, then the Gujarati for good health.
\end{quote}

\subsection*{You'll want to know: \gu{સ્વસ્થ}}
\textbf{\gu{સ્વસ્થ}} — \emph{svasth} — \textquotedblleft{}in good health; well\textquotedblright{}.

\begin{quote}
\gu{ઓપરેશન} \gu{પછી} \gu{દાદા} \gu{હવે} \gu{સ્વસ્થ} \gu{છે}.
\end{quote}

\begin{quote}
\emph{opreśan pachī dādā have svasth che.}
\end{quote}

\begin{quote}
Since the operation my grandfather is well again.
\end{quote}

\begin{cousinweb}
\textbf{\gu{સ્વસ્થ}} (\emph{svasth}) is a tatsama, Sanskrit \emph{svastha}, \textquotedblleft{}being in oneself, self-abiding\textquotedblright{}, as Monier-Williams records it: a body at ease in itself.
\end{cousinweb}

\begin{grammarlens}[title={What you've built}]
One more word for describing what you lived and felt.
\end{grammarlens}

\subsection*{Your turn}
\begin{itemize}
\raggedright
  \item \emph{Say it:} \emph{svasth}
  \item \emph{Say it:} \emph{svasth}, once more
  \item \emph{Say it:} \emph{opreśan pachī dādā have svasth che.} — Since the operation my grandfather is well again.
  \item \emph{Recall:} say the Gujarati for a diet restriction, then the Gujarati for good health, then say \emph{svasth} again
\end{itemize}

\begin{tcolorbox}[breakable,colback=teal!4,colframe=teal!35!black,title={Before you move on}]
What does \gu{સ્વસ્થ} mean? (\textquotedblleft{}In good health; well\textquotedblright{}.) Say it once more.
\end{tcolorbox}
\section[\texorpdfstring{\gu{સ્ફૂર્તિ} (sphūrti)}{સ્ફૂર્તિ (sphūrti)}]{\gu{સ્ફૂર્તિ} (sphūrti) — vigour; get-up-and-go}

\label{lesson:GU-C486-sphurti}



\begin{quote}
Before the new one: say the Gujarati for good health, then the Gujarati for in good health.
\end{quote}

\subsection*{You'll want to know: \gu{સ્ફૂર્તિ}}
\textbf{\gu{સ્ફૂર્તિ}} — \emph{sphūrti} — \textquotedblleft{}vigour; get-up-and-go\textquotedblright{}.

\begin{quote}
\gu{સવારની} \gu{કસરત} \gu{પછી} \gu{મને} \gu{સ્ફૂર્તિ} \gu{લાગે} \gu{છે}.
\end{quote}

\begin{quote}
\emph{savārnī kasrat pachī mane sphūrti lāge che.}
\end{quote}

\begin{quote}
After my morning exercise I feel full of energy.
\end{quote}

A feminine noun.

\begin{cousinweb}
\textbf{\gu{સ્ફૂર્તિ}} (\emph{sphūrti}) is a tatsama, Sanskrit \emph{sphūrti}, \textquotedblleft{}throbbing, quivering\textquotedblright{} and \textquotedblleft{}breaking forth\textquotedblright{}, as Monier-Williams records it.
\end{cousinweb}

\begin{grammarlens}[title={What you've built}]
One more word for describing what you lived and felt.
\end{grammarlens}

\subsection*{Your turn}
\begin{itemize}
\raggedright
  \item \emph{Say it:} \emph{sphūrti}
  \item \emph{Say it:} \emph{sphūrti}, once more
  \item \emph{Say it:} \emph{savārnī kasrat pachī mane sphūrti lāge che.} — After my morning exercise I feel full of energy.
  \item \emph{Recall:} say the Gujarati for good health, then the Gujarati for in good health, then say \emph{sphūrti} again
\end{itemize}

\begin{tcolorbox}[breakable,colback=teal!4,colframe=teal!35!black,title={Before you move on}]
What does \gu{સ્ફૂર્તિ} mean? (\textquotedblleft{}Vigour; get-up-and-go\textquotedblright{}.) Say it once more.
\end{tcolorbox}
\section[\texorpdfstring{\gu{ખોરાક} (khorāk)}{ખોરાક (khorāk)}]{\gu{ખોરાક} (khorāk) — food; diet}

\label{lesson:GU-C486-khorak}



\begin{quote}
Before the new one: say the Gujarati for in good health, then the Gujarati for vigour.
\end{quote}

\subsection*{You'll want to know: \gu{ખોરાક}}
\textbf{\gu{ખોરાક}} — \emph{khorāk} — \textquotedblleft{}food; diet\textquotedblright{}.

\begin{quote}
\gu{તાજો} \gu{ખોરાક} \gu{શરીર} \gu{માટે} \gu{સારો} \gu{છે}.
\end{quote}

\begin{quote}
\emph{tājo khorāk śarīr māṭe sāro che.}
\end{quote}

\begin{quote}
Fresh food is good for the body.
\end{quote}

A masculine noun.

\begin{cousinweb}
\textbf{\gu{ખોરાક}} (\emph{khorāk}) is usually traced to Persian \emph{xurāk}, \textquotedblleft{}food\textquotedblright{}, built on \emph{xurdan}, \textquotedblleft{}to eat\textquotedblright{}.
\end{cousinweb}

\begin{grammarlens}[title={What you've built}]
One more word for describing what you lived and felt.
\end{grammarlens}

\subsection*{Your turn}
\begin{itemize}
\raggedright
  \item \emph{Say it:} \emph{khorāk}
  \item \emph{Say it:} \emph{khorāk}, once more
  \item \emph{Say it:} \emph{tājo khorāk śarīr māṭe sāro che.} — Fresh food is good for the body.
  \item \emph{Recall:} say the Gujarati for in good health, then the Gujarati for vigour, then say \emph{khorāk} again
\end{itemize}

\begin{tcolorbox}[breakable,colback=teal!4,colframe=teal!35!black,title={Before you move on}]
What does \gu{ખોરાક} mean? (\textquotedblleft{}Food; diet\textquotedblright{}.) Say it once more.
\end{tcolorbox}
\section[\texorpdfstring{\gu{સુકાવું} (sukāvũ)}{સુકાવું (sukāvũ)}]{\gu{સુકાવું} (sukāvũ) — to waste away; to dry up}

\label{lesson:GU-C486-sukavu}



\begin{quote}
Before the new one: say the Gujarati for vigour, then the Gujarati for diet.
\end{quote}

\subsection*{You'll want to know: \gu{સુકાવું}}
\textbf{\gu{સુકાવું}} — \emph{sukāvũ} — \textquotedblleft{}to waste away; to dry up\textquotedblright{}.

\begin{quote}
\gu{લાંબા} \gu{તાવમાં} \gu{બાળક} \gu{સુકાઈ} \gu{ગયું}.
\end{quote}

\begin{quote}
\emph{lā̃bā tāvmā̃ bāḷak sukāī gayũ.}
\end{quote}

\begin{quote}
The child wasted away during the long fever.
\end{quote}

Past (it …): \gu{તે} \gu{છોડ} \gu{સુકાઈ} \gu{ગયો}, \emph{te choḍ sukāī gayo}.

\begin{cousinweb}
\textbf{\gu{સુકાવું}} (\emph{sukāvũ}) goes back to a reconstructed \emph{*śuṣkati}, \textquotedblleft{}dries\textquotedblright{}; CDIAL lists the Gujarati verb, \textquotedblleft{}to become emaciated\textquotedblright{} (CDIAL 12552).
\end{cousinweb}

\begin{grammarlens}[title={What you've built}]
One more word for describing what you lived and felt.
\end{grammarlens}

\subsection*{Your turn}
\begin{itemize}
\raggedright
  \item \emph{Say it:} \emph{sukāvũ}
  \item \emph{Say it:} \emph{sukāvũ}, once more
  \item \emph{Say it:} \emph{lā̃bā tāvmā̃ bāḷak sukāī gayũ.} — The child wasted away during the long fever.
  \item \emph{Recall:} say the Gujarati for vigour, then the Gujarati for diet, then say \emph{sukāvũ} again
\end{itemize}

\begin{tcolorbox}[breakable,colback=teal!4,colframe=teal!35!black,title={Before you move on}]
What does \gu{સુકાવું} mean? (\textquotedblleft{}To waste away; to dry up\textquotedblright{}.) Say it once more.
\end{tcolorbox}
\section[(dialogue)]{Looking back: from Taken Ill to Staying Well}

\label{lesson:GU-R486-review-479-486}



\begin{quote}
Say the Gujarati for unwell and paralysis.
\end{quote}

\subsection*{Your turn}
Say these aloud:

\begin{itemize}
\raggedright
  \item unwell, paralysis, a convulsion, a lump, bedridden
  \item an itchy rash, a gash, pus, to rot, to puff up
  \item to get cut, to heal, to be wounded, to ooze, to get grazed
  \item a crutch, a bump, to swell up, flat on one's back, to bang
  \item agony, a twinge, to prick, to writhe, a burning feeling
  \item a cure, a remedy, a diagnosis, to cure, an injection
  \item an operation, admitted, to put in drops, a dose, a diet restriction
  \item good health, in good health, vigour, diet, to waste away
\end{itemize}

\begin{tcolorbox}[breakable,colback=teal!4,colframe=teal!35!black,title={Before you move on}]
Unwell? (\textbf{\gu{માંદું}}, \emph{mā̃dũ}.) Paralysis? (\textbf{\gu{લકવો}}, \emph{lakvo}.) A convulsion? (\textbf{\gu{આંચકી}}, \emph{ā̃ckī}.) A lump? (\textbf{\gu{ગાંઠ}}, \emph{gā̃ṭh}.) Bedridden? (\textbf{\gu{પથારીવશ}}, \emph{pathārīvaś}.) An itchy rash? (\textbf{\gu{ખસ}}, \emph{khas}.) A gash? (\textbf{\gu{ઝખમ}}, \emph{jhakham}.) Pus? (\textbf{\gu{પરુ}}, \emph{paru}.) To rot? (\textbf{\gu{સડવું}}, \emph{saḍvũ}.) To puff up? (\textbf{\gu{ઊપસવું}}, \emph{ūpasvũ}.) To get cut? (\textbf{\gu{કપાવું}}, \emph{kapāvũ}.) To heal? (\textbf{\gu{રુઝાવું}}, \emph{rujhāvũ}.) To be wounded? (\textbf{\gu{ઘવાવું}}, \emph{ghavāvũ}.) To ooze? (\textbf{\gu{ઝરવું}}, \emph{jharvũ}.) To get grazed? (\textbf{\gu{છોલાવું}}, \emph{cholāvũ}.) A crutch? (\textbf{\gu{કાખઘોડી}}, \emph{kākhghoḍī}.) A bump? (\textbf{\gu{ઢીમડું}}, \emph{ḍhīmḍũ}.) To swell up? (\textbf{\gu{ફૂલવું}}, \emph{phūlvũ}.) Flat on one's back? (\textbf{\gu{ચત્તું}}, \emph{cattũ}.) To bang? (\textbf{\gu{અફાળવું}}, \emph{aphāḷvũ}.) Agony? (\textbf{\gu{વેદના}}, \emph{vednā}.) A twinge? (\textbf{\gu{ચસકો}}, \emph{casko}.) To prick? (\textbf{\gu{ખૂંચવું}}, \emph{khū̃cvũ}.) To writhe? (\textbf{\gu{તરફડવું}}, \emph{tarphaḍvũ}.) A burning feeling? (\textbf{\gu{દાહ}}, \emph{dāh}.) A cure? (\textbf{\gu{ઇલાજ}}, \emph{ilāj}.) A remedy? (\textbf{\gu{ઉપચાર}}, \emph{upcār}.) A diagnosis? (\textbf{\gu{નિદાન}}, \emph{nidān}.) To cure? (\textbf{\gu{મટાડવું}}, \emph{maṭāḍvũ}.) An injection? (\textbf{\gu{ઇન્જેક્શન}}, \emph{injekśan}.) An operation? (\textbf{\gu{ઓપરેશન}}, \emph{opreśan}.) Admitted? (\textbf{\gu{દાખલ}}, \emph{dākhal}.) To put in drops? (\textbf{\gu{ટપકાવવું}}, \emph{ṭapkāvvũ}.) A dose? (\textbf{\gu{માત્રા}}, \emph{mātrā}.) A diet restriction? (\textbf{\gu{પરેજી}}, \emph{parejī}.) Good health? (\textbf{\gu{આરોગ્ય}}, \emph{ārogya}.) In good health? (\textbf{\gu{સ્વસ્થ}}, \emph{svasth}.) Vigour? (\textbf{\gu{સ્ફૂર્તિ}}, \emph{sphūrti}.) Diet? (\textbf{\gu{ખોરાક}}, \emph{khorāk}.) To waste away? (\textbf{\gu{સુકાવું}}, \emph{sukāvũ}.)
\end{tcolorbox}
